\documentclass[11pt,a4paper]{article}
\usepackage[margin=22mm,headheight=15pt]{geometry}
\usepackage{fontspec}
\setmainfont{Liberation Serif}
\setsansfont{Liberation Sans}
\setmonofont{Liberation Mono}[Scale=0.86]
\usepackage{ragged2e,xurl,indentfirst}
\usepackage{ xcolor,tabularx,booktabs,enumitem,fancyhdr,titlesec,fvextra }
\usepackage[unicode,colorlinks=true,urlcolor=blue!55!black,linkcolor=blue!55!black]{hyperref}
\definecolor{ink}{HTML}{19354A}
\definecolor{accent}{HTML}{167D8D}
\titleformat{\section}{\large\sffamily\bfseries\color{ink}}{\thesection}{0.6em}{}
\titleformat{\subsection}{\normalsize\sffamily\bfseries\color{ink}}{\thesubsection}{0.6em}{}
\setlength{\parindent}{1.25cm}
\setlength{\parskip}{3pt}\justifying\setlength{\parindent}{1.25cm}\tolerance=3000\setlength{\emergencystretch}{3em}
\setlist{nosep,leftmargin=1.5em}
\renewcommand{\arraystretch}{1.12}
\pagestyle{fancy}\fancyhf{}
\fancyhead[L]{\small\sffamily ICSI405 / SEMINAR 01}
\fancyhead[R]{\small\sffamily Document RAG Chatbot}
\fancyfoot[L]{}
\fancyfoot[R]{\small\thepage}
\newcommand{\dochead}[3]{\noindent{\sffamily\small\color{accent}#1}\par\vspace{3mm}\noindent{\sffamily\bfseries\LARGE\color{ink}#2}\par\vspace{2mm}\noindent{\large #3}\par\vspace{2mm}\hrule\vspace{3mm}}
\newcommand{\note}[1]{\par\smallskip{\color{ink}\textbf{Тайлбар.} #1}\par\smallskip}
\newcommand{\pathref}[1]{\nolinkurl{#1}}
\newcommand{\src}[2]{\href{#1}{#2}}
\newcommand{\sources}{\section*{Хичээлийн эх сурвалж}\small Lecture 01, ``Why Documentation?'', слайд 4, 5, 9, 11; Seminar 01 Handout, х. 1--3. Bhatti et al., \textit{Docs for Developers}, бүлэг 1, х. 1--21; бүлэг 3, х. 58. Chinchilla, \textit{Technical Writing for Software Developers}, бүлэг 1, х. 1--8; бүлэг 2, х. 9--24. Номын дугаар нь хэвлэмэл хуудасны дугаар.\normalsize}
\begin{document}\fontsize{10.5}{13}\selectfont\titlespacing*{\section}{0pt}{10pt}{5pt}\titlespacing*{\subsection}{0pt}{7pt}{3pt}

\dochead{US-1.2 / AUDIENCE PERSONA CARD}{Зорилтот уншигч: Тэнгис}{Document RAG Chatbot}
Энэ хэсэгт өөрийгөө төслийг анх үзэж байгаа хөгжүүлэгчийн байр сууринд тавьж, ямар тайлбар хэрэгтэйг тодорхойлов.
\section{Persona Card}
\par\noindent\begin{tabularx}{\linewidth}{@{}lX@{}}\toprule
Талбар & Тодорхойлолт\\\midrule
Нэр & Тэнгис\\
Жишээ байгууллага & Nomad Learning Lab (persona-д зориулсан нэр)\\
Техникийн үүрэг & Web болон API холбож ажиллах хөгжүүлэгч\\
Уншигчийн суурь & TypeScript, React мэддэг; төслийн RAG ажиллагаатай танилцах шаардлагатай\\
Өдөр тутмын хэрэгсэл & Fedora, VS Code, Git, pnpm, DevTools, Docker Compose, curl\\
Зорилго & Local орчноос баримт оруулж, эх сурвалжтай хариуг web дээр гаргах\\
Унших хэв маяг & README-ээс эхэлж, командыг дагаж, алдааны тайлбарыг хайна\\\bottomrule
\end{tabularx}\par
\textbf{Миний тавьсан зорилго:} Нэг баримт оруулаад асуулт асуух бүх алхмыг давтаж, алдаа гарвал аль хэсгээс шалгахаа ойлгодог болох.

\textbf{Дам хэрэглэгч (end-end user):} Хөгжүүлэгч энэ баримтыг уншиж системийг ажиллуулна. Харин chatbot ашиглах хүнд хариулт нь оруулсан баримтын аль хэсгээс гарсныг харах хэрэгтэй. Түүнд embedding-ийн дотоод ажиллагааг тайлбарлахаас илүү эх сурвалжаа шалгах боломж өгөх нь чухал.
\section{Ижил төрлийн төсөлд тулгарсан гурван асуудал}
\textbf{P1. Нэвтрэлтийн cookie хадгалагдахгүй байх.} Better Auth хэрэглэгч local орчинд ажилласан app Vercel дээр refresh хийхэд нэвтрэлтээ алддаг тухай мэдээлсэн. Тэнгист зориулсан зааварт web/API хаяг, session шалгах алхмыг нэгтгэх хэрэгтэй. \src{https://stackoverflow.com/questions/79757950/better-auth-cookies-not-saving-in-frontend-after-sign-in-on-vercel-deployment-wi}{Stack Overflow: 79757950}.

\textbf{P2. Vector-ийн хэмжээс зөрөх.} SpringAI/pgvector хэрэглэгч 1596 хэмжээс хүлээж байхад 4096 ирсэн алдаа мэдээлсэн. Энэ нь өөр stack-ийн жишээ; chatbot-ийн хэмжээс гэж үзээгүй. Тэнгис model солихдоо embedder ба SQL schema-ийн нийцлийг шалгах лавлахтай байх хэрэгтэй. \src{https://stackoverflow.com/questions/78800728/springai-4096-dimensions-for-java-files-trying-to-load-in-pg-vector-db}{Stack Overflow: 78800728}.

\textbf{P3. PDF-ээс текст хоосон гарах.} pdf-parse хэрэглэгч local орчинд текст гардаг ч Azure pipeline дээр хоосон гардаг тухай мэдээлсэн. Chatbot-ийн \pathref{pdf-extractor.ts} мөн хоосон текстэд алдаа өгдөг. Текст ялгалтыг AI хариултаас салгаж оношлох, текст paste хийх өөр замыг тайлбарлах шаардлагатай. \src{https://stackoverflow.com/questions/60435500/cannot-parse-a-pdf-stream-using-pdf-parse-on-azure-pipelines-nodejs}{Stack Overflow: 60435500}.

Эдгээр асуултаас өөрийн төслийн зааварт юуг тайлбарлах хэрэгтэйг харьцуулж үзэв. Бүгд миний төсөлд гарсан алдаа биш; холбоосуудыг 2026-09-11-нд шалгасан.
\section{Өөрийн дүгнэлт}
PDF оруулсны дараа мэдээлэл олдохгүй байвал файл боловсруулагдсан эсэх, зөв баримт сонгосон эсэх, хайлт үр дүн олсон эсэхийг дарааллаар шалгах заавар хэрэгтэй гэж үзлээ.

Тиймээс эхлэх заавартаа нэг бүтэн жишээ оруулж, ready төлөв болон хариултын эх хэсгийг хаанаас харахыг тайлбарлана. Энэ дүгнэлтийг ярилцлагаас бус, төслийн ажиллагаа ба дээрх асуултуудыг харьцуулж гаргав.

\textbf{Эх сурвалж:} Seminar 01, US-1.2; Bhatti, \textit{Docs for Developers}, х. 15--17; Chinchilla, х. 7. Төслийн суурь: \pathref{AGENTS.md}, \pathref{docs/LOCAL_SETUP.md}, \pathref{apps/api/src/modules/documents/pdf-extractor.ts}.

\end{document}
